\subsubsection*{Der Chinesische Restsatz}


Ist das System simultaner Kongruenzen
\begin{align*}
z &\equiv a_1 \pmod{m_1} \\
z &\equiv a_2 \pmod{m_2} \qquad (*) \\
&\ \ \ldots \\
z &\equiv a_n \pmod{m_n}
\end{align*}
lösbar? Wenn ja, wie und wie viele Lösungen existieren?

\begin{satz}[Chinesischer Restsatz]
Sei $m, n \geq 2$. Falls $\ggT(m, n) = 1$, gilt
\[ \Zm{m \cdot n} \cong \Zm{m} \times \Zm{n} \]
\end{satz}

Sei $f : \Zm{m \cdot n} \to \Zm{m} \times \Zm{n}$, mit
\[ f(z) = (z \bmod m,\ z \bmod n) \]
Falls $m, n$ teilerfremd sind, dann ist $f$ surjektiv und injektiv, damit bijektiv, also ist $f$ ein Ring-Isomorphismus.

\begin{proof}[Bestimmung von $f^{-1}$ (Beweis des Restsatzes)]
Lösen die simultanen Kongruenzen ($\ggT(m, n) = 1$)
\begin{align*}
z &\equiv a \pmod{m} \\
z &\equiv b \pmod{n}
\end{align*}
Wenn eine Lösung $x, y \in \boldsymbol{Z}$ der Gleichung
\[ m\,x - n\,y = b - a \]
existiert, erfüllt
\[ z = m\,x + a = n\,y + b \]
beide Kongruenzen. Das ist gewährleistet (Lemma von Bezout), wenn
\[ \ggT(m, n) \mid b - a. \qedhere \]
\end{proof}

\begin{beispiel}
\begin{align*}
z &\equiv 3 \pmod{11} \\
z &\equiv 6 \pmod{8} \\
z &\equiv 1 \pmod{15}
\end{align*}
Lösen zuerst
\begin{align*}
&z \equiv 3 \pmod{11}, && a = 3 \\
&z \equiv 6 \pmod{8}, && b = 6 \\
&\text{EA mit } 11, 8, && d = b - a = 3 \\
&11 = 1 \cdot 8 + 3 && |\ + a \\
\Rightarrow\quad &z = 11 + 3 = 1 \cdot 8 + 6 \equiv 14 \pmod{88}
\end{align*}
Lösen dann
\begin{align*}
&z \equiv 14 \pmod{88}, && a = 14 \\
&z \equiv 1 \pmod{15}, && b = 6 % Hinweis: "b = 6" steht so in der Quelle; gemeint ist offenbar b = 1 (vgl. d = 1 - 14)
\\
&\text{EA mit } 88, 15, && d = b - a = 1 - 14 = -13 \\
&88 = 5 \cdot 15 + 13 && |\cdot(-1) \quad |\ + a \\
&-88 + 14 = -5 \cdot 15 + 1 \\
\Rightarrow\quad &z \equiv -88 + 14 = -74 \pmod{88 \cdot 15}
\end{align*}
Falls die kleinste Lösung $z > 0$ gesucht ist
\[ z = -74 + 1320 = 1246 \]
\end{beispiel}

\begin{bemerkung}
Der Ring-Isomorphismus $f : \Zm{m \cdot n} \to \Zm{m} \times \Zm{n}$ bildet Einheiten auf Einheiten und Nullteiler auf Nullteiler ab. Insbesondere ist die Einschränkung auf Einheiten
\[ f : \Zm{m \cdot n}^* \to \Zm{m}^* \times \Zm{n}^* \]
ein Gruppen-Isomorphismus.
\end{bemerkung}
